\chapter{Introduction}
\label{ch:intro}

\section{Background and Motivation}

An agentic AI system observes its environment, decides and acts towards a goal with little
human help. Traffic signal control fits that description closely. Every few seconds each
intersection has to choose which movements get green, using only its own detector readings,
and what one intersection does changes the traffic arriving at its neighbours. Most signals
on real roads still run fixed timing plans or simple actuated rules.

Adaptive controllers come in two families. Rule-based control such as Max-Pressure
\citep{varaiya2013max} needs no training and keeps queues bounded under stated assumptions,
but it never learns from experience. Reinforcement learning (RL) makes each intersection an
agent that learns in simulation, and on the standard benchmarks it gives lower travel times
\citep{wei2019presslight,wei2019colight}. Learned controllers also have the weaknesses that
the problem statement lists for agentic systems. Their decisions become unreliable when the
inputs drift away from the training data, they give no warning when that happens, nothing
bounds their worst case, and their decisions are hard to explain. Benchmark studies add that
reported gains are fragile and that a well-tuned independent DQN often matches newer methods
\citep{ault2021reinforcement}. This project combines the two families. A learned multi-agent
controller picks the phase, and when its own uncertainty about a decision is high, Max-Pressure
picks it instead.

\section{Problem Statement}

Problem Statement 16 asks for a study of the limits of conventional, single-agent and
multi-agent agentic frameworks on a real-world application, and for an improved system with
mathematical justification and experimental validation. Here the problem is:
\emph{given} the queue lengths, vehicle counts and current signal phase at every
intersection of a real road network (Jinan, $3\times4$; held-out city Hangzhou,
$4\times4$), \emph{choose} the next signal phase of each intersection every
$\dt$ seconds over a one-hour episode, \emph{under} minimum-green and yellow-transition
constraints and with possibly noisy or missing detector readings, \emph{so that} the
average travel time is lower than that of the baselines and the longest wait of a vehicle
stays bounded. Each action changes the inputs that follow, which makes this a sequential
decision problem.

\section{Objectives}

\begin{enumerate}
\item Compare conventional (Fixed-Time, Max-Pressure), single-agent RL (a shared DQN) and
multi-agent RL (independent DQN, IDQN) controllers on identical flows and expose their limits.
\item Design \ugmarl{}, with neighbour attention, an ensemble estimate of epistemic
uncertainty and a fallback to Max-Pressure triggered by that uncertainty, and state it
mathematically.
\item Test three hypotheses whose decision rules were fixed before the experiments:
\begin{itemize}
\item \textbf{H1 (performance):} \ugmarl{} has a lower average travel time than IDQN and
Max-Pressure on the test flow, with non-overlapping 95\,\% confidence intervals.
\item \textbf{H2 (safety):} the gate lowers the worst case (maximum per-vehicle wait, loss
of performance under detector noise) compared with the same model without the gate.
\item \textbf{H3 (coordination):} neighbour attention lowers the average travel time
compared with the same model without attention.
\end{itemize}
\item Measure robustness (detector noise, dropout, demand surge), transfer to an unseen
city, cost and explainability. The minimum acceptable outcome is a maximum wait no worse
than that of Max-Pressure.
\end{enumerate}

\section{Scope of the Work}

The work covers two grid networks with real demand, simulated in CityFlow
\citep{zhang2019cityflow} over one-hour episodes. It includes a data audit, an evaluation
protocol frozen before training, four baselines, the proposed model with three ablations,
robustness and cross-city tests, and an error analysis. Real deployment, pedestrians, public
transport, irregular networks and agents built on language models are out of scope.

\section{Contributions}

The contributions are listed in the order in which the work was done.
\begin{enumerate}
\item \textbf{Problem formulation and data design.} Signal control on a road network is
formulated as a decentralised semi-Markov decision problem with minimum-green and transition
constraints, and noisy or missing detector readings are treated as part of that problem
(\cref{sec:arch}). The conventional, single-agent and multi-agent frameworks are compared on
identical flows to measure their limits: the learned baselines beat Max-Pressure by
13--29\,s on clean data and lose to it by 345--415\,s at detector noise 0.2
(\cref{sec:comparison}). The 44 observed values are scaled by lane capacity, so no statistic
is fitted on data, and whole flow files are assigned to the splits before any other step
(\cref{sec:prep}).
\item \textbf{Mathematical modelling and theoretical analysis.} The reward, the learning
target that discounts by elapsed time, the masked ensemble loss, the uncertainty $\sigma$ and
the gate rule with its percentile threshold are written in closed form
(\cref{eq:reward,eq:target,eq:loss,eq:sigma,eq:gate}). The stability claim covers one case only.
When the gate always fires, the controller is Max-Pressure, and the queue-stability theorem of
\citet{varaiya2013max} applies under its assumptions. For the mixed policy a drift argument
is sketched without proof, and its worst case is measured instead (\cref{sec:safety}).
\item \textbf{Proposed framework.} \ugmarl{} is a multi-agent Q-learner. At every
intersection and every decision, the disagreement of an ensemble of Q-heads decides whether
the learned phase or the Max-Pressure phase is applied. Earlier ensemble Q-learners for
signals use the heads only to explore during training \citep{tan2020multiagent}, while
\ugmarl{} keeps them at execution time to choose the action. How it differs from the two
closest methods, Com-MP and the ensemble switch of Hoel et al., is set out in
\cref{sec:gaps}.
\item \textbf{Technical advancement: robustness and interpretability.} With the same weights,
the gate limits the rise in travel time at detector noise 0.2 to
\res{dtt}{ugmarl}{noise0.2}\,s, against \res{dtt}{ugmarl-nogate}{noise0.2}\,s without it.
Over six test conditions the largest regret of \ugmarl{} is \res{regret}{ugmarl}{worst}\,s,
the smallest of all controllers (\cref{tab:robust}). Each decision is logged with $\sigma$,
the gate outcome and the attention weights. CoLight \citep{wei2019colight} presents its
attention weights as an explanation. Here the weights were tested before being read that way:
they rank neighbour influence correctly, but a neighbour carries only about a seventh of the influence
of the intersection's own observation (\cref{sec:discussion}). One joint decision takes
0.81\,ms.
\item \textbf{Experimental validation and comparative analysis.} Four baselines and three
ablations run on real demand under one protocol, frozen before training, with five seeds,
95\,\% confidence intervals, a threshold sensitivity sweep and an error analysis
(\cref{ch:results}). The rules for judging the hypotheses were written before the test flow
was opened. H1 and H2 are supported; H3 is not.
\item \textbf{Generalisation and practical applicability.} The Jinan models are evaluated on
Hangzhou, a city with a different grid size, without retraining or recalibration. There
\ugmarl{} reaches \res{tt}{ugmarl}{hz-all}\,s, level with DQN and 22\,s ahead of Max-Pressure
(\cref{tab:heldout}). Detector noise, dropout and a demand surge stand in for conditions a
deployed controller would meet, and they also show a limit: under the surge the gate makes
the longest wait of a vehicle longer (\cref{sec:discussion}).
\end{enumerate}
